\paragraph{Query Embeddings.} QE models first embed entities and FOL queries into a joint low-dimensional vector space, and subsequently compute a similarity score between the \emph{entity representation} and \emph{query representation} in the latent embedding space. In general, according to the type of embedding spaces, QE-based methods can be divided into four categories: (i) geometric-based methods, \jeff{such as} GQE~\cite{hamilton2018embedding}, Q2B~\cite{query2box:ren2020}, HypE~\cite{choudhary2021self}, and ConE~\cite{zhang2021cone}, where logical queries and KG entities are embedded into a geometric vector space as points, boxes, hyperbolic, and cone shapes, respectively; (ii) distribution-based methods, such as \textsc{BetaE}~\cite{ren2020beta}, \jeff{embedding} queries to beta distributions with bounded support, and PERM~\cite{choudhary2021probabilistic}, \jeff{using} Gaussian densities to reason over KGs; 
%%%
%%
%% Paragraph Commented Out
%%%
%Compared with \textsc{BetaE}, PERM cannot handle  FOL \jeff{queries with} negation; 
%%
%%%
(iii) logic-based methods, \vicf{relating so-called set logic with FOL}
%\jeff{using} the well-studied logic to formulate the set logic in terms of FOL
~\cite{luus2021logic}; (iv) neural-based methods, \jeff{e.g.,} EMQL~\cite{sun2020faithful} \jeff{using} neural retrieval to implement logical operations. 
%%
%% Paragraph Commented Out
%%%
%but   not \jeff{supporting} logical negation. 
%%
%%
Considering QE-based methods \jeff{are} the mainstream  in the current CQA field, \textit{we mainly focus on how to construct a plug-and-play model to embed the type information for existing QE-based methods}.  

\paragraph{Other Methods.} Besides the QE-based approach, the path-based approach is another \vicf{method for} CQA, but it faces an exponential growth of the search \vicf{space} with the number of hops. For instance, CQD~\cite{arakelyan2020complex} uses a beam search method to explicitly track intermediate entities, and repeatedly combines scores from a pretrained link predictor via t-norms  to search answers while tracking intermediaries. However, CQD does not support the full set of FOL queries. %and is thus less effective. 

\paragraph{Inductive KG Completion (KGC).} In the context of KGC, there have been some works on inductive settings where test entities are not seen in the training stage. Based on the source of information used, they can be split into two categories: Using  graph structure information,  e.g.,\ subgraph or topology structures~\cite{teru2020inductive,Jiajuninductive,wang2021relational}, or  using  external information, e.g.,\ textual descriptions of entities~\cite{daza2021inductive}. However, \textit{all these methods  focus on the inductive KGC task, which can be seen as answering simpler one-hop queries.}
% , such as CraIL~\cite{teru2020inductive}, TACT~\cite{Jiajuninductive}, and PathCon~\cite{wang2021relational} exploit subgraph structures, topology structures, and relation path together with relation path between entities in an entity-independent manner, respectively. (ii) use of external information, such as BLP~\cite{daza2021inductive} and BERTRL~\cite{zha2021inductive} explore the textual descriptions of entities and pretrained language model~\cite{bert2019} to perform inductive link prediction, respectively. However, all the above mentioned methods only focus on the inductive KGC task, which is just the one-hop reasoning. To our best knowledge, there is no related work on the inductive multi-hop reasoning task over KGs.


\paragraph{Type-aware Tasks.} Type information was previously used in other tasks such as KGC or entity typing~\cite{yao2019kg,zhao2020connecting,daza2021inductive,niu2020autoeter,pan2021context}.
%For link prediction task, ~\cite{lei2020path} identifies entity types by leveraging relation type constraints to conduct knowledge graph completion. About entity typing task, ConnectE~\cite{zhao2020connecting} jointly utilizes local typing knowledge form existing entity type assertions and global triple knowledge from KGs to infer missing entity type instances. To explore the type information for any KG, AutoETER~\cite{niu2020autoeter} learns the latent type embedding of each type by regarding each relation as translation operation between the types of two entities with a relation-aware projection mechanism. ~\cite{onoe2021modeling} embed entities as d-dimensional hyperrectangles to capture hierarchies of types. CET~\cite{pan2021context} fully utilizes entities' contextual information in an independent-based and aggregated-based mechanism to infer entities' type. Although there exist above works to apply type information to the link prediction and entity typing task, 
However, these works cannot be directly used for answering FOL queries because this  requires multi-hop reasoning, producing intermediate uncertain entities.
%, and the intermediates are uncertain, which significantly complicates the task.